\documentclass[11pt]{article}

\usepackage[utf8]{inputenc}
\usepackage[T1]{fontenc}
\usepackage[margin=1in]{geometry}
\usepackage{amsmath,amssymb}
\usepackage{booktabs}
\usepackage{graphicx}
\usepackage{authblk}
\usepackage{listings}
\usepackage{xcolor}
\usepackage[numbers,sort&compress]{natbib}
\usepackage[hidelinks]{hyperref}

% --- checkmarks for the comparison table ---
\newcommand{\yes}{\checkmark}
\newcommand{\no}{$\times$}
\newcommand{\partiel}{$\sim$}

% --- code listing style ---
\definecolor{codegray}{gray}{0.95}
\lstset{
  basicstyle=\ttfamily\footnotesize,
  backgroundcolor=\color{codegray},
  keywordstyle=\color{blue!60!black},
  commentstyle=\color{green!45!black},
  stringstyle=\color{red!55!black},
  showstringspaces=false,
  breaklines=true,
  frame=single,
  language=Python,
}

\title{\texttt{manapy}: a Python framework for finite-volume methods on
unstructured meshes with unified CPU/GPU execution}

% TODO: confirm author list, order, ORCIDs and affiliations before submission
\author[1]{Imad Kissami\thanks{Corresponding author.}}
\author[1]{Ayoub Ben Hamou}
\author[1]{Mouad Haikal}
\affil[1]{Mohammed VI Polytechnic University (UM6P), Morocco}

\date{\today}

\begin{document}
\maketitle

\begin{abstract}
\texttt{manapy} is a Python framework for solving systems of conservation laws
with the cell-centered finite-volume method on unstructured, possibly hybrid,
meshes in two and three dimensions. It provides the building blocks shared by
most finite-volume solvers---mesh partitioning, halo exchange, gradient and
Laplacian reconstruction, numerical fluxes, boundary conditions, and distributed
linear solvers---together with a high-level API through which a new physical
model is assembled by reusing these operators rather than re-implementing them.
The same solver code runs on a single core, across many MPI ranks, and on CUDA
GPUs, because the numerical kernels are written against a backend abstraction
that is compiled just-in-time for the target device~\citep{lam2015numba}. Out of
the box \texttt{manapy} ships solvers for linear advection, advection--diffusion,
the inviscid Burgers equation, diffusion/Poisson/Darcy problems, the compressible
Euler equations, the shallow-water equations, and shallow-water
magnetohydrodynamics, and reads standard mesh formats through
\texttt{meshio}~\citep{meshio}.
\end{abstract}

\section{Summary}

A complete 2D advection simulation, from mesh generation to time integration and
VTK output, is set up in a few lines through the high-level API:

\begin{lstlisting}
import numpy as np
from manapy.api import Mesh, AdvectionModel

mesh = Mesh.rectangle(bounds=((0, 1), (0, 1)), n=(128, 128), cell_type="quad")
phi = mesh.field("phi",
                 init=lambda x, y, z: np.exp(-((x - 0.25)**2 + (y - 0.5)**2) / 0.01),
                 limiter="vanalbada")
model = AdvectionModel(phi, mesh, velocity=(2.0, 0.0), cfl=0.8, order=2, scheme="upwind")
model.run(T=0.25, output_every=10)      # runs serially or under mpirun -n N
\end{lstlisting}

\section{Statement of need}

Research in computational fluid dynamics and related fields routinely requires
prototyping new conservation laws, numerical fluxes, or boundary treatments and
then running them at scale. Established finite-volume codes such as
OpenFOAM~\citep{weller1998} deliver production performance but are large C++
frameworks with a steep learning curve for method development, while Python
packages such as FiPy~\citep{guyer2009fipy} and PyClaw~\citep{ketcheson2012pyclaw}
make experimentation easy but run on logically-structured grids or on the CPU
only and do not move a single implementation transparently between distributed
CPU and GPU execution on general unstructured meshes. GPU-resident frameworks
such as PyFR~\citep{witherden2014pyfr} achieve excellent performance but target
high-order flux-reconstruction schemes on specific element types rather than a
general finite-volume programming model.

\texttt{manapy} targets the gap between these: a Python-level programming model in
which a solver is expressed once, in terms of reusable finite-volume operators,
and then executed either through MPI domain decomposition on CPUs or on a GPU,
without rewriting the kernels. This is enabled by three design choices:

\begin{itemize}
  \item \textbf{Unified CPU/GPU backends.} Field data and kernels are written
  against a backend interface; the CPU backend compiles kernels with
  Numba~\citep{lam2015numba} and the GPU backend maps them to CUDA, so the
  numerics are authored once and dispatched to the active device.
  \item \textbf{A high-level, extensible, multi-physics API.} \texttt{Mesh},
  \texttt{AdvectionModel}, \texttt{DiffusionModel}, \texttt{PoissonModel} and
  \texttt{DarcyModel} let users set up a simulation in a few lines, while the
  underlying solver machinery is shared---adding a new equation reuses the
  existing flux, gradient and linear-solve operators.
  \item \textbf{Scalable parallelism.} Meshes are partitioned and coupled through
  halo exchange over \texttt{mpi4py}~\citep{dalcin2011mpi4py}; implicit problems
  are assembled and solved with distributed linear solvers via
  PETSc~\citep{petsc}, MUMPS~\citep{mumps}, Ginkgo~\citep{ginkgo}, or
  SciPy~\citep{virtanen2020scipy}.
\end{itemize}

The framework also implements advanced finite-volume ingredients---higher-order
(WENO) reconstruction, well-balanced treatments of source terms, and support for
hybrid unstructured 2D/3D grids---which makes it suitable both as a teaching and
prototyping tool and as a basis for research on new numerical methods.
\texttt{manapy} has been used to develop solvers ranging from scalar transport to
shallow-water magnetohydrodynamics, showing that the same infrastructure spans a
broad range of hyperbolic and elliptic problems.

\section{State of the field}

\begin{table}[htbp]
\centering
\small
\begin{tabular}{lccccc}
\toprule
 & \texttt{manapy} & OpenFOAM & FiPy & PyClaw & PyFR \\
\midrule
Language                    & Python & C++ & Python & Python & Python \\
Numerical method            & FV & FV & FV & FV & FR \\
General unstructured meshes & \yes & \yes & \yes & \no & \yes \\
Hybrid 2D/3D elements       & \yes & \yes & \partiel & \no & \yes \\
Distributed (MPI)           & \yes & \yes & \yes & \yes & \yes \\
GPU execution               & \yes & \no & \no & \no & \yes \\
Same solver code CPU + GPU  & \yes & \no & \no & \no & \partiel \\
High-level API for new PDEs & \yes & \partiel & \yes & \yes & \partiel \\
\bottomrule
\end{tabular}
\caption{Positioning of \texttt{manapy} among widely used open-source PDE/CFD
frameworks (FV = finite volume, FR = flux reconstruction;
\yes\ full, \partiel\ partial, \no\ none). \texttt{manapy}'s niche is a Python
finite-volume framework that combines general unstructured/hybrid meshes with a
single solver implementation that runs on both distributed CPUs and GPUs.}
\label{tab:sotf}
\end{table}

\section{Functionality}

A typical workflow reads or generates a mesh, instantiates a model, sets boundary
conditions and initial data, and advances the solution while writing VTK output
for visualization. The \texttt{manapy/examples} directory contains runnable 2D
and 3D cases for each shipped solver, and \texttt{meshes/geo} provides gmsh
\texttt{.geo} sources for the example meshes (unstructured triangular/tetrahedral,
structured quadrilateral/hexahedral, hybrid, and periodic). The public API is
exercised by the \texttt{*\_api.py} examples, while the lower-level solvers expose
the finite-volume operators directly for method development.

\section*{AI usage disclosure}

Generative-AI tools were used to assist with debugging during the development of
the software. All AI-assisted outputs were reviewed, tested, and validated by the
authors, who take full responsibility for the code and its correctness.

\section*{Acknowledgements}

% TODO: funding sources, grant numbers, contributors, and institutional support.

\bibliographystyle{plainnat}
\bibliography{paper}

\end{document}
